\documentclass[11pt,a4paper]{article}
% Common preamble for the hint sheets (\input from each file)
\usepackage[utf8]{inputenc}
\usepackage[T1]{fontenc}
\usepackage{lmodern}
\usepackage{amsmath,amssymb,amsthm,mathtools}
\usepackage[a4paper,margin=2.3cm]{geometry}
\usepackage{enumitem}
\usepackage{xcolor}
\usepackage[hidelinks]{hyperref}
\renewcommand{\leq}{\leqslant}
\renewcommand{\le}{\leqslant}
\renewcommand{\geq}{\geqslant}
\renewcommand{\ge}{\geqslant}
\setlength{\parindent}{0pt}
\setlength{\parskip}{0.35em}
% \hintstitle{Sheet title}{Class work / Homework}
\newcommand{\hintstitle}[2]{\begin{center}{\LARGE\bfseries Hints: #1\par}\vspace{0.5em}{\large #2\par}\end{center}
\noindent\textit{The hints to each problem are cumulative. Read one, try again, and open the next only if you are still stuck.}\par\medskip\hrule\medskip}
\newcommand{\problem}[1]{\par\bigskip{\large\bfseries Problem #1}\par\nopagebreak\smallskip}
\newcommand{\subpart}[1]{\par\smallskip\textbf{(#1)}\par\nopagebreak}
\newcommand{\hint}[1]{\par\smallskip\noindent\textcolor{black!70}{\textbf{Hint #1.}}\enspace}
\newcommand{\note}[1]{\par\smallskip\noindent\textit{Note.} #1\par}
\providecommand{\lt}{<}
\providecommand{\gt}{>}
\providecommand{\tightlist}{\setlength{\itemsep}{0pt}\setlength{\parskip}{0pt}}
\emergencystretch=3em

\begin{document}
\hintstitle{Methods of proof -- 6. Contradiction and proofs of impossibility}{Class work}

\problem{1}
\hint{1} Write each statement with all quantifiers explicit, then push the negation through, swapping $\forall\leftrightarrow\exists$ and negating the innermost statement. For example: ``there is a continuous $f$ on $[0,1]$ such that for every $x_0\in[0,1]$ there is $x$ with $f(x)>f(x_0)$''.
\hint{2} The others: $\exists\varepsilon>0\ \forall\delta>0\ \exists x,y$ with $|x-y|<\delta$ and $|f(x)-f(y)|\ge\varepsilon$ (for the continuity-type statement); and ``$(a_n)$ diverges'': $\forall L\in\mathbb R\ \exists\varepsilon>0\ \forall N\ \exists n\ge N$ with $|a_n-L|\ge\varepsilon$.
\hint{3} $\sqrt2$, contrapositive form: ``for all integers $p$, $q\ne0$ we have $p^2\ne2q^2$'', because the exponent of $2$ in $p^2$ is even and in $2q^2$ is odd. Descent form: if $p^2=2q^2$ with $p,q>0$, then $p=2p'$ and $q^2=2p'^2$ with $0<q<p$, a smaller solution; there is no infinite descent in $\mathbb N$.
\hint{4} The mathematics in both is the same fact about the prime $2$ (unique factorisation / parity); the contrapositive version states it directly as a valuation count, and the descent version hides it in ``$p$ even $\Rightarrow$ smaller solution''.

\problem{2}
\hint{1} Let $G$ be the set of regular elements, $|G|=5$. Check that a product of regular elements is regular and that for $a\in G$ the map $x\mapsto ax$ is injective on the whole ring. What structure does $G$ then have?
\hint{2} $G$ is a finite cancellative semigroup, hence a group. Its identity $e$ is an idempotent regular element; from $e(ex-x)=0$ and $(xe-x)e=0$ get $ex=xe=x$, so $e=1$ is the identity of the ring and all regular elements are units. Now $-1\in G$ and $(-1)^2=1$ in a group of order $5$, so $-1=1$: characteristic $2$.
\hint{3} $G$ is cyclic, generated by $g\ne1$ with $g^5=1$. The subring $\mathbb F_2[g]$ is a quotient of $\mathbb F_2[x]/(x^5-1)\cong\mathbb F_2\times\mathbb F_{16}$ (the factor $x^4+x^3+x^2+x+1$ is irreducible over $\mathbb F_2$). Since $g\ne1$, the quotient is $\mathbb F_{16}$ or the whole product; either way it has $15$ units, which are units, hence regular, in the ring. Contradiction.

\problem{3}
\hint{1} First the limit: $s_n=\sum_{k\le n}1/k^2$ is bounded, and a rational function bounded on infinitely many integers has a finite limit at $\infty$. So $D(n)=R(n)-s_n$ converges as $n\to\infty$, and since $D(n)=0$ infinitely often, $D(n)\to0$.
\hint{2} Put $Q(x)=R(x)-R(x-1)-1/x^2$, a rational function, so that $D(k)-D(k-1)=Q(k)$ and $D(n)=-\sum_{k>n}Q(k)$. If $Q\not\equiv0$, then $Q(k)$ has a constant sign for large $k$, so $D(n)\ne0$ for all large $n$: contradiction. Hence $R(x)-R(x-1)=1/x^2$ identically.
\hint{3} $R$ is not a polynomial (the difference of a polynomial is a polynomial). Let $p$ be a pole of $R$ with the largest real part: then $R(x-1)$ has a pole at $p+1$ where $R(x)$ is regular, so the left side has a pole at $p+1$, i.e.\ $p+1=0$. Similarly a pole $q$ with the smallest real part gives $q=0$. But $\operatorname{Re}q\le\operatorname{Re}p=-1$: contradiction.

\problem{4}
\hint{1} Look only at the boundary circle of the unit disc. What part of it can one disc of radius $r<1$ cover?
\hint{2} The intersection of the unit circle with a disc of radius $r<1$ is a (possibly empty) arc, any two points of which are at distance $\le2r<2$. So it contains no pair of antipodal points, hence it is a closed arc of angular length $<\pi$.
\hint{3} Two closed arcs of length $<\pi$ each have total length $<2\pi$ and cannot cover the circle.

\problem{5}
\hint{1} $\log_23=p/q$ with $p,q>0$ means $2^p=3^q$: even $=$ odd. $\sqrt2+\sqrt3=r\in\mathbb Q$ gives $r^2=5+2\sqrt6$, so $\sqrt6\in\mathbb Q$, impossible since $2$ divides $6$ exactly once (or: $\sqrt2+\sqrt3$ is a root of $x^4-10x^2+1$, which has no rational roots $\pm1$).
\hint{2} For $e=p/q$: $N=q!\bigl(e-\sum_{k\le q}1/k!\bigr)$ is an integer because $q!e=(q-1)!\,p$, and $0<N=\sum_{k>q}\frac{q!}{k!}<\frac1{q+1}+\frac1{(q+1)^2}+\dots=\frac1q\le1$.
\hint{3} Classification: integrality of $N$ uses that $e$ is a ratio $p/q$ (the denominator $q$ is cleared by $q!$), while the estimate $0<N<1$ is pure analysis; for $\log_23$ the ratio is used to produce the integer equation $2^p=3^q$. In the $\sqrt2+\sqrt3$ argument the manipulation $r\mapsto(r^2-5)/2$ uses only field operations (it works for $r$ in any field containing $\mathbb Q$, i.e.\ it shows $\mathbb Q(\sqrt2+\sqrt3)\ni\sqrt6$); only the final step $\sqrt6\notin\mathbb Q$, or the rational root test, uses the ratio of integers, and the degree argument ($[\mathbb Q(\sqrt2+\sqrt3):\mathbb Q]=4$) shows even more: it is not of degree $\le3$.

\end{document}
